\chapter*{Abstract}
\addcontentsline{toc}{chapter}{Abstract}

The disagreement between materialism and idealism is commonly framed as a dispute over what reality is fundamentally made of, yet both positions begin by dividing reality through conceptual vocabularies whose adequacy is rarely examined before the ontological argument begins. This work asks a prior question: which descriptive language preserves the structure of reality with the least loss, and which mathematics is required when ordinary language or lower-order formal descriptions collapse distinctions the phenomena themselves make? Mathematics is treated neither as an infallible mirror of reality nor as a merely convenient notation, but as a family of formal languages whose state spaces, transformations, metrics, quotients, symmetries, and invariant algebras determine what can be distinguished and what is identified.

The mathematical spine begins with a relational state and separates state identity from form identity. Two reduced states may be genuinely different while instantiating the same form, so persistence need not require stasis. Coupled oscillators and Stuart--Landau systems make the distinction between collective motion and internal relational variation explicit. Finite-group projectors on the sphere provide exact measures of Platonic symmetry coherence, while character theory separates mathematical admissibility from dynamical selection. In the real $\ell=5$ representation of $SO(3)$, the quartic reduced energy is shown to possess an exact five-dimensional local minimum manifold at an explicit regular point; three tangent directions are rotational and two are genuine shape-flat directions. Sixth-order invariant structure contains four genuinely new invariant directions, and an explicit sextic invariant has positive-definite intrinsic Hessian on both unresolved quartic shape directions. This establishes a concrete case in which a richer descriptive algebra distinguishes structure unavailable to a lower-order one.

Complex wave conjugation supplies a complementary result. Conjugation preserves norm, density, symmetry coherence, and stabilizer while reversing directed current, demonstrating that fine-state transformation can coexist with coarse-form identity. Platonic duality and complex conjugation are nevertheless distinct in general, and a fixed geometry can support inequivalent dynamical laws. The work therefore treats Platonic form and wave dynamics as mathematical archetypes of invariance and transformation rather than as literal stand-ins for idealism and materialism.

The second half of the book separates descriptive adequacy from ontological sufficiency. A better language can preserve more structure without becoming the reality it describes. The physical--experiential question is therefore tested through a separate evidential hierarchy. Observable phenomenological evidence is represented by a latent variable $Z_\varphi$ that is explicitly not identified with experience itself. Increasingly comprehensive physical descriptions are tested for sufficiency using conditional information, partial information decomposition, dynamical closure, intervention, cross-context transfer, and complexity-controlled model comparison. The final distinction is between mathematical admissibility and phenomenal occupancy: a formalism may have room for consciousness without showing that consciousness occupies that structure. The program is explicitly defeasible. If physical enrichment absorbs every phenomenologically relevant residual, relational physicalism gains support; if independently anchored phenomenological structure remains necessary after intervention and complexity-controlled comparison, relational dual-aspect realism gains empirical motivation; if the alternatives remain observationally equivalent, underdetermination is accepted as the result.
